\begin{frame}{자주 묻는 질문 (4.3)}
  \small
  \begin{itemize}
    \item \textbf{Q. k-means가 항상 수렴하나요?}
      그렇다 --- 매 반복마다 SSE가 절대 늘지 않고 $0$ 이상이므로
      \textbf{유한한 시간 안에} 수렴이 보장. 다만
      \textbf{어디로} 수렴하는지(전역 vs 나쁜 지역최적해)는
      보장되지 않음 = ``자주 하는 실수''의 핵심
    \item \textbf{Q. $k$를 몰라도 쓸 수 있나요?}
      k-means 자체는 $k$를 \textbf{미리} 정해야 함.
      $k$를 데이터로부터 자동으로 원하면: 엘보우로 여러 $k$를
      비교하거나, 4.4절 EM/GMM처럼 각 클러스터의 ``적합도''를
      \textbf{확률적으로} 비교할 수 있는 모델로 넘어가기
    \item \textbf{Q. 왜 항상 ``원형'' 덩어리만 찾나요?}
      \textbf{구조적으로} 결정됨 --- 두 클러스터의 경계는
      ``두 중심점까지 거리가 같은 점들의 집합''이고, 2차원에선
      \textbf{직선}(수직이등분선) $\Rightarrow$ 각 클러스터 영역은
      반드시 \textbf{볼록} 다각형(보로노이 셀). 반달·고리 모양
      경계를 그릴 \textbf{자유 자체가 모델 안에 없음} ---
      파라미터 튜닝으로 극복할 수 없는 \textbf{구조적} 한계
    \item \textbf{Q. 클러스터가 비어(멤버 0)질 수 있나요?}
      이론적으로 가능 --- 어떤 중심점도 어떤 점에게도 ``가장
      가까운''이 아닌 상황. \texttt{if cluster else
      centroids[i]} 분기가 바로 이걸 막음(빈 클러스터의
      ``평균''은 정의되지 않음 $\to$ 중심점 유지가 관례)
    \item \textbf{Q. 차원이 높으면 k-means도 안 쓰나요?}
      4.2절의 차원의 저주가 \textbf{그대로} 적용 --- ``가장
      가까운 중심점'' 구분이 점점 무뎌짐. 실무 표준:
      \textbf{Week 14 PCA로 차원 축소}한 뒤 2$\sim$3차원 공간에서
      클러스터링
  \end{itemize}
\end{frame}
